% Module 1 section file. Included by ../module1_stellarator_equilibrium_geometry.tex.
\section{Coordinate transformation and interpolation}
\label{sec:transform_interpolation}

\subsection{Why transformations are unavoidable}
Different calculations prefer different representations:
\begin{itemize}[leftmargin=1.6em]
\item coil and wall geometry are naturally expressed in Cartesian or cylindrical coordinates;
\item equilibrium solvers use internal flux coordinates and spectral coefficients;
\item guiding-center solvers often prefer Boozer coordinates or direct Cartesian fields;
\item wave solvers require metric coefficients and parallel derivatives;
\item transport solvers frequently use a one-dimensional normalized flux coordinate.
\end{itemize}
Module 1 must therefore define transformations, not only raw field values.

\subsection{Forward and inverse mappings}
The forward mapping
\begin{equation}
(s,\theta,\zeta)\mapsto\mathbf x
\end{equation}
is evaluated from the surface representation. The inverse mapping
\begin{equation}
\mathbf x\mapsto(s,\theta,\zeta)
\end{equation}
is generally nonlinear and may require root finding. The inverse can be nonunique outside the nested-surface region or near self-intersections, so its domain must be restricted.

A Newton inverse solves
\begin{equation}
\mathbf r(s,\theta,\zeta)
=\mathbf x(s,\theta,\zeta)-\mathbf x_{\mathrm{target}}
=\mathbf 0.
\end{equation}
The update uses the coordinate-basis matrix
\begin{equation}
\begin{bmatrix}
\partial_s\mathbf x & \partial_\theta\mathbf x & \partial_\zeta\mathbf x
\end{bmatrix}.
\end{equation}
Near a coordinate singularity or a poorly shaped surface, this matrix can be ill-conditioned.

\subsection{Interpolation requirements}
A downstream solver may evaluate geometry at points not present on the equilibrium grid. Interpolation should preserve:
\begin{itemize}[leftmargin=1.6em]
\item periodicity in angular coordinates;
\item continuity across field-period boundaries;
\item regularity at the magnetic axis;
\item stated order of accuracy;
\item consistency between a field and its derivatives.
\end{itemize}
Interpolating $B$, $\nabla B$, and $\boldsymbol\kappa$ independently can violate derivative identities. A differentiable spectral representation is preferable when practical.

\subsection{Conservative remapping of flux functions}
When a profile such as density or pressure is transferred from one radial grid to another, pointwise interpolation does not necessarily conserve integrated particle number or energy. A conservative remap should preserve an integral such as
\begin{equation}
\mathcal N
=\int_0^1 n(s)V'(s)\,ds.
\label{eq:particle_integral}
\end{equation}
For cell averages, this is accomplished by matching integrals over overlapping radial cells rather than simply matching values at cell centers.

